\documentclass[11pt]{article}
\usepackage{amsmath,amssymb,booktabs}
\usepackage[margin=1in]{geometry}
\title{Step 204: Unnormalized \(\Xi_{\mathrm{matrix\_source}}\) Attempt}
\date{2026-05-16}
\begin{document}
\maketitle

\section*{Verdict}
\[
\boxed{\texttt{V\_unnormalized\_partial}.}
\]

\section*{Inherited Matrix Context}

Step 162 declares \(P_\eta\) as the projection onto the closed pulled
zero-evaluator span \(H_\eta\), and \(C_\ell=(I-P_\infty)M_{m_\ell}P_\infty\).
Step 177 records
\[
c_{ij}(\ell)=\langle e_i,C_\ell e_j\rangle,\qquad
e_i=\eta_{\rho_i,0}^{1/2}/\|\eta_{\rho_i,0}^{1/2}\|,
\]
on \(Rho_{\rm fin}=\{\rho_1,\rho_2,\rho_3\}\), \(a_0=1/2\), \(K_{\rm fin}=\{0\}\).

\section*{Unnormalized Trace Identity}

Let \(\kappa_i=\kappa_{1/2,\rho_i,0}\) and
\[
G_{ij}=\langle \kappa_i,\kappa_j\rangle.
\]
Then the finite projection is
\[
P_\eta x=\sum_{i,j}\kappa_i(G^{-1})_{ij}\langle \kappa_j,x\rangle.
\]
For the full Hilbert--Schmidt residual,
\[
\|C_\ell P_\eta\|_{\rm HS}^2
=\operatorname{tr}_{H_\eta}(C_\ell^*C_\ell)
=\operatorname{tr}(G^{-1}H),
\qquad
H_{ij}=\langle C_\ell\kappa_i,C_\ell\kappa_j\rangle.
\]
Thus the full residual requires \(H\), not only
\[
c_{ij}=\langle\kappa_i,C_\ell\kappa_j\rangle.
\]
The \(c\)-only identity
\[
\operatorname{tr}(G^{-1}c^*G^{-1}c)
\]
is the Hilbert--Schmidt norm of the compressed operator
\(P_\eta C_\ell P_\eta\).  It equals the full residual only with an additional
range-in-\(H_\eta\) hypothesis, or if the finite diagnostic is explicitly
defined as the compression.

\section*{Gram Matrix}

The Step 204 off-diagonal convention uses
\[
G_{ij}=K_{1/2}^{\Gamma}(\rho_j,\overline{\rho_i}),\qquad i\ne j,
\]
with the Step 203 L'Hopital diagonal.  Representative entries are
\[
G_{11}=6.661579628286\cdot10^{-10},\quad
G_{12}=-1.913153364500\cdot10^{-13},\quad
G_{22}=1.287139952474\cdot10^{-14},\quad
G_{33}=2.394335172196\cdot10^{-17}.
\]
The nominal determinant is
\[
\det(G)\approx 2.10033807151\cdot10^{-40},
\]
and \(\|G^{-1}\|_F\approx 4.09992478506\cdot10^{16}\).  The diagonal error
bounds inherited from Step 203 still dominate the diagonal entries.

\section*{\(c\)-Matrix Attempt}

The unnormalized entries are
\[
c_{ij}=\langle\kappa_i,(I-P_\infty)M_{m_\ell}P_\infty\kappa_j\rangle.
\]
The unnormalized formulation removes division by \(\sqrt{G_{ii}}\), but it
does not remove the need for concrete boundary-line samples
\[
\kappa_i(\tau)=T_{1/2}^*K_{1/2}^{\Gamma}(\cdot,\rho_i)(1/2+i\tau)
\]
and the full \(K_\infty^{op}\) PSWF quadrature.  Burnol's \(E\)-kernel gives
pairings and the projected reproducing kernel, but this step did not obtain a
certified numerical \(T_a^*\)-sampling formula.

\section*{Conclusion}

\(\Xi_{\mathrm{matrix\_source}}\) is not decided.  The next gate is either:
derive the full \(H_{ij}=\langle C_\ell\kappa_i,C_\ell\kappa_j\rangle\) matrix
for the true Hilbert--Schmidt residual, or explicitly demote the finite
diagnostic to the compressed \(P_\eta C_\ell P_\eta\) norm and compute
\(c_{ij}\) with certified \(\kappa\)-boundary samples.

\end{document}

